% =============================================================
%  PREAMBLE — Common Lisp: Guida Pratica
%  Compilare con: pdflatex main  (x2)
% =============================================================

% --- Encoding e lingua ---
\usepackage[T1]{fontenc}
\usepackage[utf8]{inputenc}
\usepackage[italian]{babel}

% --- Font ---
\usepackage{lmodern}
\usepackage{microtype}

% --- Layout pagina ---
\usepackage[
  a4paper,
  top=2.5cm, bottom=2.8cm,
  inner=2.8cm, outer=2.2cm,
  marginparwidth=1.5cm
]{geometry}

% --- Colori ---
\usepackage[dvipsnames,svgnames,table]{xcolor}

\definecolor{clBg}{RGB}{248,248,248}       % sfondo listing
\definecolor{clKeyword}{RGB}{0,0,180}      % keyword Lisp
\definecolor{clComment}{RGB}{100,100,100}  % commenti
\definecolor{clString}{RGB}{0,120,0}       % stringhe

% Colori difficoltà
\definecolor{colOne}{RGB}{130,130,130}     % ★      grigio
\definecolor{colTwo}{RGB}{34,110,34}       % ★★     verde
\definecolor{colThree}{RGB}{30,80,160}     % ★★★    blu
\definecolor{colFour}{RGB}{200,100,0}      % ★★★★   arancio
\definecolor{colFive}{RGB}{180,30,30}      % ★★★★★  rosso

% Colori box
\definecolor{boxEssFrame}{RGB}{30,80,160}
\definecolor{boxEssBack}{RGB}{235,241,252}
\definecolor{boxColFrame}{RGB}{110,110,110}
\definecolor{boxColBack}{RGB}{245,245,245}
\definecolor{boxHintFrame}{RGB}{180,130,0}
\definecolor{boxHintBack}{RGB}{255,252,220}
\definecolor{boxSolFrame}{RGB}{34,110,34}
\definecolor{boxSolBack}{RGB}{236,248,236}
\definecolor{boxRagFrame}{RGB}{140,50,140}
\definecolor{boxRagBack}{RGB}{248,236,248}

% --- Matematica ---
\usepackage{amsmath}
\usepackage{amssymb}

% --- Listings (codice Lisp) ---
\usepackage{listings}

\lstdefinelanguage{CommonLisp}{%
  sensitive=false,
  morecomment=[l]{;},
  morestring=[b]",
  morekeywords=[1]{%
    defun, defmacro, defparameter, defvar, defconstant,
    defstruct, defclass, defmethod, defgeneric,
    let, let*, flet, labels, setf, setq,
    if, cond, when, unless, case, ecase,
    and, or, not,
    progn, prog1, prog2,
    loop, dotimes, dolist, do, do*,
    lambda, function,
    funcall, apply,
    mapcar, maplist, mapcan, mapc,
    remove-if, remove-if-not, find-if, find-if-not,
    every, some, notany, notevery,
    reduce, count-if, count,
    sort, stable-sort,
    return, return-from,
    values, multiple-value-bind, multiple-value-list,
    destructuring-bind,
    handler-case, handler-bind, error, warn, signal,
    with-open-file, open, close,
    format, print, princ, prin1, terpri, fresh-line,
    read, read-line, read-char,
    car, cdr, cons, list, append, reverse, length,
    first, second, third, fourth, fifth,
    rest, last, nth, nthcdr, butlast,
    null, atom, listp, consp, numberp, stringp,
    symbolp, functionp, integerp, floatp, characterp,
    eq, eql, equal, equalp,
    quote, backquote,
    make-hash-table, gethash, remhash, maphash, hash-table-count,
    make-array, aref, array-dimensions,
    make-instance, slot-value, initialize-instance,
    in-package, defpackage, use-package, export, import,
    compile, load, require, provide,
    t, nil%
  },
  morekeywords=[2]{%
    +, -, *, /, mod, rem, expt, sqrt, abs,
    floor, ceiling, truncate, round,
    min, max, 1+, 1-,
    =, /=, <, >, <=, >=,
    char, char-code, code-char, char-upcase, char-downcase,
    string, string-upcase, string-downcase, string-length,
    coerce, type-of, typep, check-type%
  }
}

\lstdefinestyle{lispstyle}{%
  language=CommonLisp,
  basicstyle=\ttfamily\small,
  keywordstyle=[1]{\color{clKeyword}\bfseries},
  keywordstyle=[2]{\color{clKeyword}},
  commentstyle=\itshape\color{clComment},
  stringstyle=\color{clString},
  backgroundcolor=\color{clBg},
  frame=single,
  framesep=4pt,
  framerule=0.4pt,
  rulecolor=\color{gray!50},
  xleftmargin=8pt,
  xrightmargin=4pt,
  breaklines=true,
  breakatwhitespace=true,
  showstringspaces=false,
  tabsize=2,
  numbers=none,
  upquote=true,
  extendedchars=true,
  literate=
    {→}{{$\rightarrow$}}1
    {★}{{\textcolor{colOne}{$\bigstar$}}}1
    {à}{{\`a}}1  {è}{{\`e}}1  {é}{{\'{e}}}1  {ì}{{\`{\i}}}1
    {ò}{{\`o}}1  {ù}{{\`u}}1
    {À}{{\`A}}1  {È}{{\`E}}1  {É}{{\'{E}}}1
    {Ì}{{\`I}}1  {Ò}{{\`O}}1  {Ù}{{\`U}}1
}

\lstset{style=lispstyle}

% Shorthand per listing inline
\newcommand{\cl}[1]{\lstinline[style=lispstyle]{#1}}

% --- tcolorbox ---
\usepackage[most]{tcolorbox}
\tcbuselibrary{listings, breakable, skins}

% Stelle di difficoltà
\newcommand{\unaStella}{\textcolor{colOne}{$\bigstar$}}
\newcommand{\dueStelle}{\textcolor{colTwo}{$\bigstar\bigstar$}}
\newcommand{\treStelle}{\textcolor{colThree}{$\bigstar\bigstar\bigstar$}}
\newcommand{\quattroStelle}{\textcolor{colFour}{$\bigstar\bigstar\bigstar\bigstar$}}
\newcommand{\cinqueStelle}{\textcolor{colFive}{$\bigstar\bigstar\bigstar\bigstar\bigstar$}}

% Mappa numero di stelle → colore frame
\newcommand{\starcolor}[1]{%
  \ifcase#1
    \color{colOne}%   0 (non usato)
  \or\color{colOne}%  1
  \or\color{colTwo}%  2
  \or\color{colThree}% 3
  \or\color{colFour}% 4
  \else\color{colFive}% 5
  \fi
}

% ---- Box ESERCIZIO ESSENZIALE ◆ ----
% Uso: \begin{essbox}{n}{titolo}{stelle} ... \end{essbox}
\newenvironment{essbox}[3]{%
  \medskip
  \begin{tcolorbox}[
    enhanced,
    breakable,
    colframe=boxEssFrame,
    colback=boxEssBack,
    fonttitle=\bfseries\small,
    title={{\small\textcolor{boxEssFrame}{$\blacklozenge$}}\enspace Esercizio~#1\enspace---\enspace#2\enspace
      \hfill\ifcase#3\or\unaStella\or\dueStelle\or\treStelle\or\quattroStelle\or\cinqueStelle\fi},
    arc=3pt,
    boxrule=1.2pt,
    left=6pt, right=6pt, top=4pt, bottom=4pt
  ]
}{%
  \end{tcolorbox}
  \medskip
}

% ---- Box ESERCIZIO COLLATERALE ○ ----
% Uso: \begin{colbox}{n}{titolo}{stelle} ... \end{colbox}
\newenvironment{colbox}[3]{%
  \medskip
  \begin{tcolorbox}[
    enhanced,
    breakable,
    colframe=boxColFrame,
    colback=boxColBack,
    fonttitle=\bfseries\small,
    title={{\small\textcolor{boxColFrame}{$\circ$}}\enspace Esercizio~#1\enspace---\enspace#2\enspace
      \hfill\ifcase#3\or\unaStella\or\dueStelle\or\treStelle\or\quattroStelle\or\cinqueStelle\fi},
    arc=3pt,
    boxrule=0.8pt,
    borderline={0.8pt}{0pt}{boxColFrame, dashed},
    left=6pt, right=6pt, top=4pt, bottom=4pt
  ]
}{%
  \end{tcolorbox}
  \medskip
}

% ---- Box SUGGERIMENTO ----
\newenvironment{hintbox}[1]{%
  \medskip
  \begin{tcolorbox}[
    enhanced,
    breakable,
    colframe=boxHintFrame,
    colback=boxHintBack,
    fonttitle=\bfseries\small,
    title={Suggerimento --- Esercizio~#1},
    arc=3pt,
    boxrule=0.8pt,
    left=6pt, right=6pt, top=4pt, bottom=4pt
  ]
  \small
}{%
  \end{tcolorbox}
  \medskip
}

% ---- Box SOLUZIONE ----
\newenvironment{solbox}[1]{%
  \medskip
  \begin{tcolorbox}[
    enhanced,
    breakable,
    colframe=boxSolFrame,
    colback=boxSolBack,
    fonttitle=\bfseries\small,
    title={Soluzione --- Esercizio~#1},
    arc=3pt,
    boxrule=0.8pt,
    left=6pt, right=6pt, top=4pt, bottom=4pt
  ]
  \small
}{%
  \end{tcolorbox}
  \medskip
}

% ---- Box RAGIONAMENTO (per ★★★★/★★★★★) ----
\newenvironment{ragionamentoBox}[1]{%
  \medskip
  \begin{tcolorbox}[
    enhanced,
    breakable,
    colframe=boxRagFrame,
    colback=boxRagBack,
    fonttitle=\bfseries\small,
    title={Ragionamento --- Esercizio~#1},
    arc=3pt,
    boxrule=0.8pt,
    left=6pt, right=6pt, top=4pt, bottom=4pt
  ]
  \small
}{%
  \end{tcolorbox}
  \medskip
}

% Alias brevi per compatibilita' (comandi -> ambienti)
% Non usare questi per listing, usare direttamente gli ambienti
\newcommand{\essenzialeBox}[4]{\begin{essbox}{#1}{#2}{#3}#4\end{essbox}}
\newcommand{\collateraleBox}[4]{\begin{colbox}{#1}{#2}{#3}#4\end{colbox}}
\newcommand{\hintBox}[2]{\begin{hintbox}{#1}#2\end{hintbox}}
\newcommand{\soluzioneBox}[2]{\begin{solbox}{#1}#2\end{solbox}}

% --- Intestazioni capitolo ---
\usepackage{titlesec}
\titleformat{\chapter}[display]
  {\normalfont\Large\bfseries\color{boxEssFrame}}
  {\chaptertitlename\ \thechapter}{10pt}{\Huge}
\titlespacing*{\chapter}{0pt}{20pt}{30pt}

% --- Header/Footer ---
\usepackage{fancyhdr}
\pagestyle{fancy}
\fancyhf{}
\fancyhead[LE,RO]{\small\thepage}
\fancyhead[LO]{\small\nouppercase{\rightmark}}
\fancyhead[RE]{\small\nouppercase{\leftmark}}
\renewcommand{\headrulewidth}{0.4pt}
\fancypagestyle{plain}{%
  \fancyhf{}
  \fancyfoot[C]{\small\thepage}
  \renewcommand{\headrulewidth}{0pt}
}

% --- Hyperref (deve essere quasi l'ultimo) ---
\usepackage[
  colorlinks=true,
  linkcolor=boxEssFrame,
  urlcolor=colTwo,
  citecolor=colThree,
  bookmarksnumbered=true,
  bookmarksopen=true,
  pdftitle={Common Lisp --- Guida Pratica},
  pdfauthor={},
  pdfsubject={Common Lisp},
  pdfkeywords={Common Lisp, programmazione funzionale, esercizi}
]{hyperref}

% --- Multicol per cheat sheet ---
\usepackage{multicol}
\usepackage{booktabs}
\usepackage{array}
\usepackage{enumitem}

% Impaginazione: evita vedove e orfane
\widowpenalty=10000
\clubpenalty=10000

% Etichette esercizi, hint, soluzioni
% Uso: \label{ex:42}  \label{hint:42}  \label{sol:42}

% Comandi utili
\newcommand{\kw}[1]{\texttt{\bfseries #1}}  % keyword inline
\newcommand{\fn}[1]{\texttt{#1}}            % funzione inline
\newcommand{\result}{{\color{clComment}; $\rightarrow$\ }}

% Nomi simboli
\newcommand{\essential}{{\color{boxEssFrame}$\blacklozenge$}}
\newcommand{\optional}{{\color{boxColFrame}$\circ$}}
